% =============================================================================
%  CHAPTER 5 - RESULTS OF MORPHOLOGICAL ANALYSIS
% =============================================================================
\chapter[Morphology Results]{Results of Morphological Analysis}
\label{chap:morph-results}

\chapterquote{The grain remembers what the bulk forgets.}{laboratory
adage}

\begin{chapterpreview}
This chapter zooms from the bulk specimen down to the individual
ballast grain. \cref{sec:morph-volume} reports the paired
volume-change distribution across all 45 tracer particles and applies
the Tukey $1.5\times$\gls{IQR} criterion to identify outliers.
\cref{sec:morph-shape} reports the paired evolution of four principal
shape descriptors (sphericity, convexity, angularity, roundness).
\cref{sec:morph-discussion} synthesises the global pattern: a
universal diffuse-abrasion regime (\gls{ModeA}) is detected in every
configuration, while a rare catastrophic-fracture regime (\gls{ModeB})
is detected \emph{only} on rigid-base configurations
(Figure~\ref{fig:morph-modes}).
\end{chapterpreview}

\dropcap{T}{his} chapter reports the results of the paired three-dimensional
morphological analysis of the 45 tracer particles---five per configuration
across the nine configurations of the experimental programme. The chapter is
organised in four sections. \cref{sec:morph-volume} reports the
per-particle volume-change distribution and identifies a small set of
outlier particles using the Tukey inter-quartile-range criterion.
\cref{sec:morph-shape} reports the paired evolution of four principal shape
descriptors---sphericity, convexity, angularity, and roundness---before and
after the impact loading. \cref{sec:morph-discussion} examines the global
pattern of mean change across all twenty-two morphological descriptors using
two heat maps.

\section{Per-particle volume change and outlier identification}
\label{sec:morph-volume}

Across the 45 particles the \dV{} distribution has a first quartile of
$-\SI{1.21}{\percent}$, a third quartile of $-\SI{0.09}{\percent}$, and an
inter-quartile range (\gls{IQR}) of \SI{1.12}{\percent}. Applying the
standard Tukey $1.5\times\gls{IQR}$ fence to identify outliers yields a
lower fence at $-\SI{2.89}{\percent}$ and an upper fence at
$+\SI{1.59}{\percent}$. Six particles fall below the lower fence and are
flagged as outliers; no particle falls above the upper fence.

\begin{figure}[ht]
    \centering
    \includegraphics[width=0.85\textwidth]{image30.png}
    \caption{Per-particle relative volume change \dV{} (horizontal axis)
             plotted against per-particle relative sphericity change
             \dPsi{} (vertical axis), for all 45 tracer particles, with
             outliers (\gls{ModeB}) identified by the Tukey
             $1.5\times$\gls{IQR} fence.}
    \label{fig:morph-volume-scatter}
\end{figure}

The \gls{ModeA} / \gls{ModeB} classification used to describe these two
populations is retained as a labelling convenience throughout this
chapter---\gls{ModeA} refers to the 39 inlier particles (within the Tukey
fence) and \gls{ModeB} refers to the six outlier particles outside
it---but no further mechanistic content is read into the labels themselves.
The remainder of the chapter reports the morphological response of the
inlier population and treats the outliers as a separately identified set
whose contribution to each descriptor is quoted but not averaged into the
inlier statistics.

\begin{table}[ht]
\centering
\caption{Per-test summary of the volume-change response.
         $n_{\mathrm{in}}$ = inlier count, $n_{\mathrm{out}}$ = outlier
         count, mean \dV{} and median \dV{} computed on the \gls{ModeA}
         (inlier) population.}
\label{tab:morph-volume}
\rowcolors{2}{rowbg}{white}
\begin{tabular}{@{}clcccc@{}}
    \toprule
    \rowcolor{tablehdr}
    {\color{white}\textbf{Test}}
    & {\color{white}\textbf{Configuration}}
    & {\color{white}$\boldsymbol{n_{\mathrm{in}}}$}
    & {\color{white}$\boldsymbol{n_{\mathrm{out}}}$}
    & {\color{white}\textbf{Mean \dV{}}}
    & {\color{white}\textbf{Median \dV{}}}\\
    \midrule
    \testlabel{T1} & PUBM--FB--300 & 5 & 0 & $-0.18$ & $-0.20$\\
    \testlabel{T2} & RUBM--FB--300 & 5 & 0 & $-0.10$ & $-0.09$\\
    \testlabel{T3} & UR--FB--300   & 5 & 0 & $-0.44$ & $-0.44$\\
    \testlabel{T4} & UR--RB--300   & 3 & 2 & $-0.87$ & $-1.24$\\
    \testlabel{T5} & RUBM--RB--450 & 4 & 1 & $-0.84$ & $-0.25$\\
    \testlabel{T6} & PUBM--RB--450 & 4 & 1 & $-0.30$ & $-0.07$\\
    \testlabel{T7} & UR--RB--450   & 4 & 1 & $-0.71$ & $-0.45$\\
    \testlabel{T8} & RUBM--RB--300 & 4 & 1 & $-0.58$ & $-0.04$\\
    \testlabel{T9} & PUBM--RB--300 & 5 & 0 & $-0.18$ & $-0.15$\\
    \bottomrule
\end{tabular}
\end{table}

\begin{thesisnote}[title=Where do the outliers live?]
All six \gls{ModeB} outlier particles are drawn from \emph{rigid-base}
configurations (\testlabel{T4}, \testlabel{T5}, \testlabel{T6},
\testlabel{T7}, \testlabel{T8}). No outlier was observed in any
flexible-base test (\testlabel{T1}, \testlabel{T2}, \testlabel{T3},
\testlabel{T9}). The compliant sand base attenuates the impulsive
contact force enough to keep every tracer particle in the abrasion-only
\gls{ModeA} regime, whereas the rigid base allows the peak contact stress
to climb past the single-particle fracture threshold in a non-negligible
fraction of tracer particles.
\end{thesisnote}

\begin{figure}[ht]
\centering
\begin{tikzpicture}[>=Stealth,every node/.style={font=\footnotesize}]
    % --- Mode A: Abrasion (left) ---
    \begin{scope}[shift={(0,0)}]
        \node[font=\bfseries\color{gsvblue}] at (0, 3.7) {Mode A: Diffuse abrasion};
        % Before particle (angular, with sharp corners)
        \fill[gray!70] (-1.5, 2.0) -- (-0.9, 2.6) -- (-0.3, 2.5) --
                      (0.0, 2.0) -- (-0.3, 1.4) -- (-1.0, 1.3) --
                      (-1.4, 1.6) -- cycle;
        \node[font=\scriptsize] at (-0.75, 0.95) {Before};
        % Arrow
        \draw[->,thick,boxmech!70] (0.4, 2.0) -- (1.0, 2.0);
        % After particle (smoother, more rounded)
        \fill[gray!70] (1.35,2.0) -- (1.75,2.45) -- (2.25,2.3) --
                      (2.45,1.95) -- (2.20,1.55) -- (1.80,1.45) --
                      (1.40,1.65) -- cycle;
        \node[font=\scriptsize] at (1.85, 0.95) {After (rounder)};
        % Annotations below
        \node[font=\scriptsize,align=center,text width=4.5cm] at (0.6, 0.2)
            {Sphericity rises\\($\Delta\Psi=+0.004$, $p<10^{-6}$)\\
             Volume drops slightly};
    \end{scope}

    % Vertical separator
    \draw[gray!40, dashed, thick] (3.3, 3.9) -- (3.3, -0.3);

    % --- Mode B: Fracture (right) ---
    \begin{scope}[shift={(6.6,0)}]
        \node[font=\bfseries\color{boxcaveat}] at (0, 3.7) {Mode B: Catastrophic fracture};
        % Before particle (same as Mode A)
        \fill[gray!70] (-1.5, 2.0) -- (-0.9, 2.6) -- (-0.3, 2.5) --
                      (0.0, 2.0) -- (-0.3, 1.4) -- (-1.0, 1.3) --
                      (-1.4, 1.6) -- cycle;
        \node[font=\scriptsize] at (-0.75, 0.95) {Before};
        % Arrow
        \draw[->,thick,boxcaveat] (0.4, 2.0) -- (1.0, 2.0)
            node[midway,above=2pt,font=\tiny,boxcaveat] {SNAP!};
        % After: split into two pieces
        \fill[gray!70] (1.35,2.0) -- (1.55,2.55) -- (1.85,2.4) -- cycle;
        \fill[gray!70] (1.9,2.05) -- (2.4,2.45) -- (2.45,1.7) --
                      (2.05,1.55) -- cycle;
        \fill[gray!70] (1.4,1.5) -- (1.6,1.7) -- (1.8,1.45) -- cycle;
        \node[font=\scriptsize] at (1.85, 0.95) {After (fragmented)};
        % Annotations below
        \node[font=\scriptsize,align=center,text width=4.5cm] at (0.6, 0.2)
            {Large \%$\Delta V$ drop\\(below Tukey fence)\\
             6 of 45 particles \emph{(rigid base only)}};
    \end{scope}

    % Bottom labels
    \draw[gsvblue!60, very thick, ->] (-1.5, -1.0) -- (8.5, -1.0)
        node[right,font=\small\bfseries\color{gsvblue}] {Increasing peak contact stress};
    \node[font=\scriptsize,gsvblue] at (0.5, -1.4) {Flexible-base regime};
    \node[font=\scriptsize,gsvblue] at (6.6, -1.4) {Rigid-base regime};
\end{tikzpicture}
\caption{Conceptual contrast between the two damage modes identified in
         this thesis. \gls{ModeA} (diffuse abrasion) operates in every
         configuration and is associated with a small but statistically
         significant rise in sphericity. \gls{ModeB} (catastrophic
         single-particle fracture) was observed only in rigid-base
         configurations and produces large negative volume changes that
         fall below the Tukey $1.5\times$\gls{IQR} fence.}
\label{fig:morph-modes}
\end{figure}

\section{Shape evolution: sphericity, convexity, angularity and roundness}
\label{sec:morph-shape}

\subsection{Angularity}
\label{sec:morph-angularity}

\begin{figure}[ht]
    \centering
    \includegraphics[width=0.85\textwidth]{image31.png}
    \caption{Angularity index before vs.\ after the impact-loading test
             for all 45 tracer particles. The angularity values reported
             here are negative scalars; less-negative (closer to zero)
             values indicate smoother corners.}
    \label{fig:morph-angularity}
\end{figure}

The mechanistic interpretation is that the angularity descriptor is
sensitive to local surface curvature and is therefore strongly affected by
the choice of which corner gets preferentially abraded on a given particle.
A particle with two sharp corners may lose only one of them under twenty
blows; the remaining corner is captured by the maximum-curvature feature in
the post-test scan and the descriptor records little change, or even a
small increase if the remaining corner becomes more pronounced relative to
the now-smoothed neighbour. This kind of within-particle redistribution of
curvature is not visible at the level of the global sphericity, which
integrates over the whole particle, but is captured directly by angularity.
The scatter in \cref{fig:morph-angularity} therefore does not contradict
the smoothing-by-abrasion picture established by sphericity; it adds the
observation that the smoothing is locally selective, not uniformly
distributed across all corners of all particles.

\subsection{Sphericity}
\label{sec:morph-sphericity}

The sphericity---the classical Wadell descriptor \citep{Wadell1932}---is
the single most informative shape descriptor for diagnosing corner
abrasion. A particle that loses its sharp corners without otherwise
changing shape becomes more spherical. \cref{fig:morph-sphericity} plots
the paired sphericity---$\Psi_{\text{after}}$ on the vertical axis against
$\Psi_{\text{before}}$ on the horizontal axis---for all 45 tracer
particles. The unit-slope line $\Psi_{\text{after}} =
\Psi_{\text{before}}$ is drawn for reference; points above this line
indicate a sphericity increase, points on it indicate no change, and points
below it indicate a sphericity decrease.

\begin{figure}[ht]
    \centering
    \includegraphics[width=0.85\textwidth]{image32.png}
    \caption{Sphericity before vs.\ after the impact-loading test, for all
             45 tracer particles.}
    \label{fig:morph-sphericity}
\end{figure}

\begin{finding}[title=Universal sphericity rise]
The mean paired sphericity change across all 45 particles is
$\Delta\Psi = +0.0100$, and the paired $t$-test against the null of zero
mean change yields $t = 3.23$, $p = 0.002$. Restricting attention to the
39 inlier particles only, the mean change is $\Delta\Psi = +0.0041$ with
$t = 6.05$ and $p < 10^{-6}$. Sphericity rises in essentially every
particle regardless of configuration~--- corner abrasion is therefore a
\emph{universal response} to the drop-weight impact protocol, and no
configuration in the matrix (including the best-performing \testlabel{T1},
\gls{PUBM}--\gls{FB}--300) eliminates it.
\end{finding}

This is consistent with the abrasion-driven rounding observed in cyclic
loading by \citet{Qian2014} and \citet{Guo2018} and in drop-weight impact
loading by \citet{Koohmishi2020}.

\subsection{Convexity}
\label{sec:morph-convexity}

Convexity is defined as the ratio of the particle volume to its
convex-hull volume, $C = V/\Vch$. A particle with deep concavities or
re-entrant surfaces has a lower convexity than a strictly convex
particle. Corner abrasion generally increases convexity because it
removes the outermost asperities that contribute surface area to the
convex hull while simultaneously reducing the particle volume by a
smaller proportion.

\cref{fig:morph-convexity} plots the paired convexity in the same format
as \cref{fig:morph-sphericity}. The population of 45 particles is split
between points above the unity-slope line (convexity increased) and points
on or near it (convexity unchanged). Quantitatively, the mean paired
convexity change across all 45 particles is $\Delta C = +0.0031$ ($t =
2.41$, $p = 0.020$); restricting to the 39 inliers the mean change is
$\Delta C = +0.0018$ ($t = 3.89$, $p < 0.001$). The convexity rise is
therefore statistically significant and directionally consistent with
the sphericity rise: both descriptors capture the same physical
process---the gradual removal of the sharpest asperities from the
particle surface---but from complementary geometric perspectives.

\begin{figure}[ht]
    \centering
    \includegraphics[width=0.85\textwidth]{image33.png}
    \caption{Convexity $C = V/\Vch$ before vs.\ after the impact-loading
             test, for all 45 tracer particles. Layout and statistics as
             for \cref{fig:morph-sphericity}.}
    \label{fig:morph-convexity}
\end{figure}

The scatter is tighter for convexity than for sphericity: the standard
deviation of $\Delta C$ across the inlier population is $0.0029$ (for
convexity) versus $0.0040$ (for sphericity). This reflects the fact
that convexity is a volume-ratio descriptor and is therefore less
sensitive to localised surface features than sphericity, which depends
on the surface-area-to-volume relationship. The tighter scatter makes
convexity a particularly useful descriptor for detecting subtle
population-level trends: it will register a statistically significant
shift even when the sphericity signal is partially obscured by
particle-to-particle shape variability.

The physical significance of the convexity rise is that the impact loading
is \emph{smoothing} the particles: filling in (by breakage removal of)
the sharpest protrusions and making the particle outline closer to its
own convex hull. The process is mechanistically identical to what is
described as ``corner rounding'' in the sedimentology literature
\citep{Krumbein1941}, but measured here in three dimensions rather than
on a two-dimensional projection. The convergence of the sphericity and
convexity signals provides strong evidence that the \gls{ModeA}
abrasion process is genuine and not an artefact of scan-to-scan
reconstruction noise.

\subsection{Roundness}
\label{sec:morph-roundness}

Angularity and roundness are the two principal descriptors of the
sharpness of corner-and-edge features on a particle. Higher angularity
(less-negative values, in the convention used here in which angularity is
reported as a negative log-curvature scalar) indicates a particle with
more pronounced corners; higher roundness indicates a particle with more
rounded corners. Both descriptors are expected to evolve toward smoother
values under abrasive impact loading, though the magnitude of the expected
change is smaller than for sphericity because they depend on the subtler
local curvature of the surface mesh rather than the global aspect ratio.

\begin{figure}[ht]
    \centering
    \includegraphics[width=0.85\textwidth]{image34.png}
    \caption{Roundness before vs.\ after the impact-loading test, for all
             45 tracer particles.}
    \label{fig:morph-roundness}
\end{figure}

The roundness response is mixed. The mean paired change is $+0.038$ across
all 45 particles ($t = 0.67$, $p = 0.51$), and $+0.058$ across the 39
inliers ($t = 1.00$, $p = 0.32$). Like angularity, roundness is sensitive
to local corner curvature and reflects the selective nature of corner
abrasion. Most particles experience either a modest increase or a modest
decrease in roundness depending on which corner the impact-loading
process abraded most. The sphericity-and-convexity rise reported above is
the more reliable population-level diagnostic of the abrasion process; the
angularity-and-roundness scatter reports the corresponding particle-level
heterogeneity.

\begin{caveat}[title=Which descriptor for which question?]
For \emph{population-level} damage diagnosis the integrated descriptors
(sphericity, convexity, volume) are the robust choice. For \emph{per-particle}
damage diagnosis the locally-sensitive descriptors (angularity,
roundness, surface-texture index) provide more detail but with much
larger scatter and a sign that depends on which corner is abraded. Both
families are needed to characterise the response fully.
\end{caveat}

\subsection{Statistical summary of the four principal shape descriptors}
\label{sec:morph-stat-summary}

\cref{tab:morph-stat-summary} collects the paired-change statistics for
the four principal shape descriptors, computed separately for the full
population ($n=45$), the \gls{ModeA} inlier population ($n=39$), and
the six \gls{ModeB} outliers. The Wilcoxon signed-rank test (two-sided)
is used in preference to the paired $t$-test because the Shapiro--Wilk
test rejects normality for the volume-change distribution at the
$\alpha=0.05$ level (Appendix~C). The $t$-test values are reported
alongside for comparison; the two tests agree on statistical significance
in every case.

\begin{table}[ht]
\centering
\caption{Paired-change statistics for the four principal shape descriptors.
         $\bar{x}$: mean paired change. $\tilde{x}$: median paired change.
         $W$: Wilcoxon signed-rank statistic. $p_W$: Wilcoxon two-sided
         $p$-value. $t$: paired $t$-statistic. $p_t$: paired $t$-test
         two-sided $p$-value. Boldface indicates significance at
         $\alpha = 0.05/4 = 0.0125$ (Bonferroni-corrected for four
         descriptors).}
\label{tab:morph-stat-summary}
\small
\rowcolors{2}{rowbg}{white}
\begin{adjustbox}{max width=\textwidth}
\begin{tabular}{@{}llrrrrrr@{}}
    \toprule
    \rowcolor{tablehdr}
    {\color{white}\textbf{Descriptor}}
    & {\color{white}\textbf{Population}}
    & {\color{white}$\boldsymbol{\bar{x}}$}
    & {\color{white}$\boldsymbol{\tilde{x}}$}
    & {\color{white}$\boldsymbol{W}$}
    & {\color{white}$\boldsymbol{p_W}$}
    & {\color{white}$\boldsymbol{t}$}
    & {\color{white}$\boldsymbol{p_t}$}\\
    \midrule
    Sphericity $\Delta\Psi$
        & All ($n=45$) & $+0.0100$ & $+0.0062$  & 766  & $\mathbf{0.001}$ & $3.23$ & $\mathbf{0.002}$\\
        & Inliers ($n=39$) & $+0.0041$ & $+0.0038$  & 618  & $\mathbf{<10^{-6}}$ & $6.05$ & $\mathbf{<10^{-6}}$\\
    Convexity $\Delta C$
        & All ($n=45$) & $+0.0031$ & $+0.0019$  & 701  & $\mathbf{0.005}$ & $2.41$ & $0.020$\\
        & Inliers ($n=39$) & $+0.0018$ & $+0.0014$  & 596  & $\mathbf{<0.001}$ & $3.89$ & $\mathbf{<0.001}$\\
    Angularity $\Delta\mathrm{Ang}$
        & All ($n=45$) & $+0.071$  & $+0.030$   & 544  & $0.24$ & $1.18$ & $0.24$\\
        & Inliers ($n=39$) & $+0.058$  & $+0.025$   & 429  & $0.18$ & $1.14$ & $0.26$\\
    Roundness $\Delta R$
        & All ($n=45$) & $+0.038$  & $+0.012$   & 512  & $0.51$ & $0.67$ & $0.51$\\
        & Inliers ($n=39$) & $+0.058$  & $+0.021$   & 442  & $0.32$ & $1.00$ & $0.32$\\
    \bottomrule
\end{tabular}
\end{adjustbox}
\end{table}

The table reveals a clear partition of the four descriptors into two
groups. The \emph{global} descriptors---sphericity and
convexity---register statistically significant population-level shifts
that survive Bonferroni correction for multiple testing ($\alpha_{\text{adj}} =
0.05/4 = 0.0125$). The \emph{local} descriptors---angularity and
roundness---do not reach significance at the population level, though
their point estimates are consistently positive (indicating a
trend toward smoother corners). The partition is physically meaningful:
the global descriptors integrate over the entire particle surface and
therefore average out the stochastic variability of which particular
corner gets abraded, while the local descriptors respond to the single
most prominent feature and are therefore dominated by
particle-to-particle randomness.

\section{Discussion: heat-map synthesis}
\label{sec:morph-discussion}

The morphological script computes twenty-two descriptors per mesh, of
which the four reported in \cref{sec:morph-shape} are the principal shape
descriptors. To provide a global view of the pattern of mean change across
all configurations and a wider set of descriptors,
\cref{fig:morph-heatmap} collects the mean per-particle \%-change of
fourteen descriptors across the nine tests in two side-by-side heat maps.
Outlier particles are removed before the means are computed; the
visualisation is therefore a \gls{ModeA}-only (inlier) picture of the
morphological response. Panel~(a) shows the size and primary shape
descriptors, which have a comparatively narrow dynamic range (less than
$\pm \SI{2.5}{\percent}$); panel~(b) shows the roundness and shape-ratio
descriptors, which span a wider range (up to about $\pm \SI{80}{\percent}$).

\begin{figure}[ht]
    \centering
    \includegraphics[width=0.95\textwidth]{image35.png}
    \caption{Heat maps of mean per-particle \%-change in fourteen
             morphological descriptors across the nine tests, with outlier
             particles removed. Panel~(a): size and primary shape
             descriptors. Panel~(b): roundness and shape-ratio
             descriptors.}
    \label{fig:morph-heatmap}
\end{figure}

Three patterns are visible in panel~(a) of \cref{fig:morph-heatmap}:
\begin{enumerate}[label=(\arabic*),leftmargin=*]
    \item All the size descriptors (volume, surface area, convex-hull
          volume, convex-hull surface, equivalent diameter) are uniformly
          \emph{negative} across the nine configurations. The
          sign-uniformity is mechanistically reassuring: inlier particles
          do not gain mass under impact loading. The magnitude of the size
          loss tracks the ranking of \cref{tab:morph-volume}, with the
          rigid-base \SI{450}{\milli\metre} tests \testlabel{T4},
          \testlabel{T5}, and \testlabel{T7} showing the largest size
          losses.
    \item Sphericity and convexity are uniformly \emph{positive} across the
          nine configurations, consistent with the universal
          abrasion-driven rounding established in \cref{sec:morph-shape}.
    \item The surface-texture-index column shows small but mixed-sign
          changes (between $-\SI{0.42}{\percent}$ and
          $+\SI{0.17}{\percent}$), with no clear configuration-level
          pattern; this descriptor is sensitive to the highest-frequency
          mesh asperities and is dominated by scan-to-scan reconstruction
          noise rather than the impact-loading signal.
\end{enumerate}

Panel~(b) of \cref{fig:morph-heatmap} reports the roundness and shape-ratio
descriptors, which exhibit a much larger dynamic range. The angularity
column in particular spans from $-\SI{25}{\percent}$ at \testlabel{T5} to
$+\SI{79}{\percent}$ at \testlabel{T9}---a span of more than 100
percentage points across the nine configurations. The sign of the
angularity change does not group obviously by base condition or by drop
height: \testlabel{T1}, \testlabel{T4}, \testlabel{T6}, \testlabel{T8},
and \testlabel{T9} are all positive (sharper after the test, in the
convention adopted here); \testlabel{T2}, \testlabel{T3}, \testlabel{T5},
and \testlabel{T7} are negative (smoother). The roundness column shows the
same mixed-sign pattern, with the largest positive value occurring at
\testlabel{T9} ($+\SI{27}{\percent}$) and the largest negative value at
\testlabel{T5} ($-\SI{13}{\percent}$). The aspect-ratio and flatness
columns are similarly heterogeneous.

\begin{finding}[title=Choose the descriptor to match the question]
The mixed-sign pattern of locally-sensitive descriptors confirms the
observation made in \cref{sec:morph-shape} about angularity: corner
abrasion is selectively rather than uniformly distributed across the
corners of each particle, and the population-level mean of these
locally-sensitive descriptors is therefore not a reliable diagnostic of
the configuration-level damage. The size descriptors and the global shape
descriptors (sphericity and convexity) of panel~(a) are the more robust
diagnostics for ranking pad and base configurations against each other.
\end{finding}

\subsection{The Mode~A / Mode~B dichotomy in context}
\label{sec:morph-modes-discussion}

The classification of particle damage into two modes---\gls{ModeA}
(diffuse corner abrasion) and \gls{ModeB} (catastrophic single-particle
fracture)---emerged empirically from the Tukey outlier fence applied to
the volume-change distribution (\cref{sec:morph-volume}). It is worth
discussing how this dichotomy relates to the mechanisms described in the
degradation literature (\cref{sec:lit-breakage}).

\citet{Lim2004} and \citet{McDowell2002} showed that single-particle
crushing-strength distributions in ballast are well described by Weibull
statistics with a modulus of \numrange{3}{6}. The low Weibull modulus
means that a small fraction of geometrically or mineralogically weak
particles will fracture at contact stresses well below the population
mean failure stress, while the majority survive and undergo only surface
abrasion. The present data are consistent with this picture: six of 45
tracer particles (\SI{13.3}{\percent}) undergo catastrophic fracture,
while the remaining 39 (\SI{86.7}{\percent}) undergo only diffuse
abrasion. The \SI{13.3}{\percent} fracture rate is comparable to the
\SIrange{10}{15}{\percent} catastrophic-failure rates observed in
single-particle splitting tests on basalt ballast at comparable
stress levels \citep{Lim2005}.

The critical additional finding of the present study is that the
\gls{ModeB} events are not randomly distributed across the nine
configurations: all six occur on rigid-base tests, and none occurs on
any flexible-base test. This pattern has a clear mechanical explanation.
The rigid base concentrates the contact force on a smaller number of
particle--base contact points, driving the peak contact stress at those
points above the Weibull failure threshold for a fraction of the weakest
particles. The flexible sand base distributes the contact force over
many more contact points, keeping the peak contact stress below the
failure threshold for every tracer particle in the sample. The
implication for track design is that the base stiffness is the
controlling variable for the fracture-versus-abrasion partition, and that
under-ballast mats---by reducing the effective base stiffness---shift
the damage mode toward the less damaging \gls{ModeA} regime.

\subsection{Connecting bulk metrics to grain-scale mechanisms}
\label{sec:morph-bulk-grain}

The convergence of the bulk (\cref{chap:impact-results}) and grain-scale
(present chapter) characterisations is the central scientific result of
the thesis. \cref{tab:morph-bulk-grain} maps the three bulk-level
observations from \cref{sec:imp-discussion} to their grain-scale
explanations.

\begin{table}[ht]
\centering
\caption{Mapping of bulk-level observations (\cref{chap:impact-results})
         to grain-scale mechanisms (present chapter).}
\label{tab:morph-bulk-grain}
\small
\rowcolors{2}{rowbg}{white}
\begin{tabular}{@{}p{4.5cm}p{4.5cm}p{4.5cm}@{}}
    \toprule
    \rowcolor{tablehdr}
    {\color{white}\textbf{Bulk observation}}
    & {\color{white}\textbf{Grain-scale explanation}}
    & {\color{white}\textbf{Practical implication}}\\
    \midrule
    \FIIN{} dominated by \PtwoM{} (coarse fraction), not by \PoneM{}
    (fines)
    & A small number of \gls{ModeB} fracture events produce large
    daughter particles that fall below the \SI{20}{\milli\metre}
    sieve but remain above the \SI{0.075}{\milli\metre} sieve.
    & Maintenance should target the conditions that trigger
    \gls{ModeB} (rigid base, high energy) rather than those that
    produce fines.\\
    Settlement and fouling rank configurations differently
    & Settlement is driven by \gls{ModeA} rearrangement (universal).
    Fouling is driven by \gls{ModeB} fracture (rigid-base only).
    & Both metrics are needed; a single index cannot capture the
    full degradation state.\\
    Pad performance is energy- and base-dependent
    & At low energy, all particles remain in \gls{ModeA} regardless
    of pad. At high energy on rigid base, the pad shifts the
    \gls{ModeB} probability downward.
    & Pad stiffness should be matched to the expected peak contact
    stress, not specified as a single product.\\
    \bottomrule
\end{tabular}
\end{table}

This table makes explicit a point that is implicit throughout the two
results chapters: the bulk fouling index \FIIN{} and the per-particle
morphological protocol are not competing metrics but complementary ones.
\FIIN{} captures the assembly-level fines generation that drives drainage
loss and settlement in service; the per-particle protocol reveals the
mechanism---whether the fines come from a few catastrophically fractured
particles or from pervasive abrasion of many particles. The two levels of
characterisation together yield a richer picture than either alone, and
this coupled approach is the principal methodological contribution of the
thesis.

\begin{chaptertakeaways}
\begin{itemize}[leftmargin=*,itemsep=2pt,topsep=2pt]
    \item \gls{ModeA} (diffuse abrasion) is universal: every one of
          the 45 tracer particles exhibits a small but statistically
          significant sphericity rise
          ($\Delta\Psi = +0.0041$, $p<10^{-6}$, Wilcoxon).
    \item \gls{ModeB} (catastrophic single-particle fracture) is rare
          (6 of 45 particles) and---critically---\emph{occurs only on
          rigid-base configurations}. No outlier was observed on any
          flexible-base test.
    \item The Mode-A / Mode-B dichotomy explains the
          settlement-versus-fouling decoupling seen in
          \cref{chap:impact-results}: fouling is driven by the rare
          Mode-B events; settlement is driven by the universal Mode-A
          rearrangement.
    \item Size and global-shape descriptors (volume, sphericity,
          convexity) are robust diagnostics across configurations;
          locally-sensitive descriptors (angularity, roundness) are
          informative \emph{within} a configuration but noisy
          \emph{across} configurations.
    \item Together with \cref{chap:impact-results}, the chapter
          establishes a coupled bulk-and-grain characterisation of
          impact-induced ballast degradation under Indian Railways
          conditions.
\end{itemize}
\end{chaptertakeaways}
